\documentclass[11pt,reqno]{amsart}

\usepackage[T1]{fontenc}
\usepackage[utf8]{inputenc}
\usepackage{amsmath,amssymb,amsthm}
\usepackage[letterpaper,margin=0.75in]{geometry}
\usepackage{booktabs}
\usepackage[hidelinks]{hyperref}

\DeclareMathOperator{\vol}{vol}
\DeclareMathOperator{\conv}{conv}
\DeclareMathOperator{\area}{area}
\DeclareMathOperator{\spn}{span}

%% Artifact references: a second reference group pointing into the repository
%% rather than the literature, numbered [R1], [R2], ... so that it cannot be
%% confused with the bibliography. Set \repobase and \repocommit once; every
%% link and every citation follows.
\newcommand{\repobase}{https://github.com/csabamedgyeswork/inscribed-polytopes}
\newcommand{\repocommit}{main}
\newcommand{\repohref}{\href{\repobase}{\texttt{\repobase}}}
% \nolinkurl rather than \texttt so a long path may break after a slash.
\newcommand{\artpath}[1]{\href{\repobase/tree/\repocommit/#1}{\nolinkurl{#1}}}

\newcounter{artifact}
\newenvironment{artifactlist}
  {\sloppy % entries are mostly path, so prefer stretched spaces to overrun
   \begin{list}{[R\arabic{artifact}]}
     {\usecounter{artifact}%
      \setlength{\labelwidth}{3.1em}\setlength{\leftmargin}{3.7em}%
      \setlength{\labelsep}{0.6em}\setlength{\itemsep}{0.15ex}%
      \setlength{\parsep}{0pt}}}
  {\end{list}}
\newcommand{\artifact}[3]{\item\label{art:#1}\artpath{#2}, #3}
\newcommand{\artref}[1]{\hyperref[art:#1]{[R\ref*{art:#1}]}}

\newtheorem{theorem}{Theorem}[section]
\newtheorem{proposition}[theorem]{Proposition}
\newtheorem{lemma}[theorem]{Lemma}
\newtheorem{corollary}[theorem]{Corollary}

\theoremstyle{definition}
\newtheorem{definition}[theorem]{Definition}
\newtheorem{remark}[theorem]{Remark}
\newtheorem*{notation}{Notation}

% A quoted statement kept under the name and number of its source.
\newcommand{\citedname}{}
\newtheorem*{citedstatement}{\citedname}
\newenvironment{cited}[2]
  {\renewcommand{\citedname}{#1}\begin{citedstatement}[#2]}
  {\end{citedstatement}}

\title[The inscribed $4$-polytope of maximal volume with nine vertices]
      {The inscribed $4$-polytope of maximal volume\\ with nine vertices}
\author{Csaba Medgyes}
\email{csaba.medgyes.work@gmail.com}

\begin{document}

\begin{abstract}
  We prove that a $4$-polytope with nine vertices inscribed in the unit sphere
  $\mathbb{S}^{3}$ has volume at most $\frac{5}{24}\sqrt{10 + 2\sqrt{5}}$, and
  that, up to orthogonal transformations, equality holds only for the convex
  hull of a square and a regular pentagon inscribed in $\mathbb{S}^{3}$ in
  orthogonal $2$-planes through the origin. For the combinatorial type of this
  polytope, whose boundary complex is the join of a $4$-cycle and a $5$-cycle,
  we write the signed volume as the pairing,
  through the Hodge star, of the area vectors of the two cycles, and bound each
  area vector by finite Fourier analysis; the bound holds for every choice of
  nine points on the sphere. Since an optimal polytope is
  simplicial, by a result of Horv\'ath and L\'angi, it remains to exclude the
  other simplicial types, which we do by sum-of-squares certificates checked in
  exact arithmetic.
\end{abstract}

\maketitle

\section{Introduction}
\label{sec:intro}

Which convex polytope with $n$ vertices, inscribed in the unit sphere
$\mathbb{S}^{d-1}$ of $\mathbb{R}^{d}$, has the largest volume? Fejes T\'{o}th
raised this question in three dimensions \cite{fejes-toth-1964}, and it is
listed among the open problems of discrete geometry
\cite{brass-moser-pach-2005}. There, besides the classical cases $n = 4, 6,
12$ (see \cite[Section~1.2]{hoehner-ledford-2026}), Berman and Hanes
determined the optimal polyhedra with at most eight vertices
\cite{berman-hanes-1970}, and Hoehner and Ledford recently those with nine
\cite{hoehner-ledford-2026}. For more vertices only numerical results are
known \cite{mutoh-2003}, \cite{hardin-sloane-smith}.

In higher dimensions even less is known. Horv\'{a}th and L\'{a}ngi extended
the necessary condition for optimality of Berman and Hanes to every dimension,
and with it showed that optimal polytopes are simplicial
\cite[Lemma~1]{horvath-langi-2016}. They determined the optimal polytopes with
$n = d + 1$ vertices, the regular simplices
\cite[Corollary~1]{horvath-langi-2016}; with $n = d + 2$, the convex hulls of
two regular simplices of dimensions $\lfloor d/2\rfloor$ and $\lceil
d/2\rceil$ in orthogonal subspaces \cite[Theorem~2]{horvath-langi-2016}; and,
for odd $d$, with $n = d + 3$ \cite[Theorem~3]{horvath-langi-2016}. For even
$d$ their argument for $n = d + 3$ excludes cyclic polytopes, and they asked
whether an optimal polytope can be cyclic
\cite[Problem~2]{horvath-langi-2016}.

In dimension four these results settle $n = 5$ and $n = 6$. In earlier work
we answered the question of Horv\'{a}th and L\'{a}ngi for $d = 4$, which
settles $n = 7$, and showed that the regular cross-polytope is optimal for
$n = 8$ \cite{medgyes-maxvol-d4}. The optima for $n = 6$, $7$ and $8$ share a
structure: each is the convex hull of two regular polygons with $m$ and $l$
vertices, inscribed in $\mathbb{S}^{3}$ in orthogonal $2$-planes through the
origin, with $(m, l) = (3, 3)$, $(3, 4)$ and $(4, 4)$. To our knowledge, no case with $d = 4$
and $n \geq 9$ has been settled.

This paper treats the first open case in dimension four, $n = 9$, where the
pattern suggests a square and a regular pentagon. The method of
\cite{medgyes-maxvol-d4} bounds the volume on each combinatorial type of
simplicial polytope by a sum-of-squares certificate, checked in exact
arithmetic. Two things change with nine vertices. The number of types grows
from $37$ \cite{grunbaum-sreedharan-1967} to $1142$
\cite{altshuler-bokowski-steinberg-1980}. And the volume of the square--pentagon pair is
irrational, while a certificate gives a rational bound, so on the type of the
pair no certificate is sharp.

We prove that the square--pentagon pair is, up to orthogonal transformations,
the unique optimum (Theorem~\ref{thm:main}). Certificates bound the volume on
every other type strictly below that of the pair. On the type of
the pair, whose boundary complex is the join of a $4$-cycle and a $5$-cycle,
we prove a sharp
inequality instead. Its two ingredients are a factorisation of the signed
volume through the Hodge star, and a bound on the area vector of a closed
polygon with unit vertices, proved by expanding the vertex coordinates in the
real discrete Fourier basis, as Schoenberg did for plane polygons
\cite{schoenberg-1950}.

Section~\ref{sec:prelim} collects the definitions and known results we use.
Section~\ref{sec:lemmas} proves three lemmas about the join of two cycles,
among them the factorisation of the signed volume, and
Section~\ref{sec:polygon} the inequality for closed polygons.
Section~\ref{sec:pair} proves the sharp bound within the join type,
Section~\ref{sec:rivals} bounds all other types, and
Section~\ref{sec:main} combines the two.

\section{Preliminaries}
\label{sec:prelim}

\subsection{Inscribed polytopes}

\begin{cited}{Notation}{{\cite[Section~1]{horvath-langi-2016}}}
  Throughout the paper, by a polytope we mean a convex polytope. The vertex set
  of a polytope $P$ is denoted by $V(P)$. We denote the family of
  $d$-dimensional polytopes, with $n$ vertices and inscribed in the unit sphere
  $\mathbb{S}^{d-1}$, by $\mathcal{P}_{d}(n)$. We denote $d$-dimensional
  volume by $\vol_{d}$, and set
  $v_{d}(n) = \max\{\vol_{d}(P) : P \in \mathcal{P}_{d}(n)\}$.
\end{cited}

\begin{remark}[{\cite[Remark~2.1]{medgyes-maxvol-d4}}]
  \label{rem:attained}
  For $n \geq d + 1$, $v_{d}(n)$ is attained by some
  $P \in \mathcal{P}_{d}(n)$.
\end{remark}

\begin{definition}[{\cite[Definition~2.1, Section~2.2 and Proposition~2.16]{ziegler-1995}}]
  A \emph{facet} of a $d$-polytope is a face of dimension $d-1$, and a
  \emph{ridge} is a face of dimension $d-2$. A polytope is \emph{simplicial} if
  every facet is a simplex. Two polytopes are \emph{combinatorially
  equivalent} if their face lattices are isomorphic; an equivalence class is a
  \emph{combinatorial type}.
\end{definition}

Equivalently, two polytopes are combinatorially equivalent if their
facet--vertex incidence matrices differ only by permutations of rows and
columns \cite[Section~15.1.1]{henk-richter-gebert-ziegler-2017}. For
simplicial polytopes, which is the form in which we use the definition, this
means that a bijection between the vertex sets carries the vertex sets of the
facets of one onto those of the other.

\begin{definition}[{\cite[Definition~1]{horvath-langi-2016}}]
  \label{def:propertyz}
  Let $P \in \mathcal{P}_{d}(n)$ be a $d$-polytope with
  $V(P) = \{p_{1}, p_{2}, \ldots, p_{n}\}$. If for each $i$, there is an open
  set $U_{i} \subset \mathbb{S}^{d-1}$ such that $p_{i} \in U_{i}$, and for any
  $q \in U_{i}$, we have
  \begin{equation*}
    \vol_{d}\bigl(\conv((V(P) \setminus \{p_{i}\}) \cup \{q\})\bigr)
      \leq \vol_{d}(P),
  \end{equation*}
  then we say that $P$ satisfies Property Z.
\end{definition}

\begin{remark}[{\cite{berman-hanes-1970}, \cite[Remark~2.4]{medgyes-maxvol-d4}}]
  \label{rem:propertyz}
  Every $P \in \mathcal{P}_{d}(n)$ with $\vol_{d}(P) = v_{d}(n)$ satisfies
  Property Z.
\end{remark}

\begin{lemma}[{\cite[Lemma~1]{horvath-langi-2016}}]
  \label{lem:simplicial}
  Every $P \in \mathcal{P}_{d}(n)$ satisfying Property Z is simplicial.
\end{lemma}

\subsection{Signed volume and certificates}

For $v_{1}, \ldots, v_{d} \in \mathbb{R}^{d}$, $\det(v_{1}, \ldots, v_{d})$
denotes the determinant of the matrix with columns $v_{1}, \ldots, v_{d}$.

\begin{definition}[{\cite[Definition~3.1]{medgyes-maxvol-d4}}]
  \label{def:signed-volume}
  An \emph{orientation} of a finite set is an ordering of its elements up to
  even permutations. Let $\mathcal{C} = (n, \mathcal{F})$, where
  $\mathcal{F}$ is a finite set of ordered $d$-tuples
  $F = (F_{1}, \ldots, F_{d})$ of distinct elements of $\{1, \ldots, n\}$,
  called facets. Deleting $F_{a}$ gives a ridge, with the orientation of
  $(F_{1}, \ldots, \widehat{F_{a}}, \ldots, F_{d})$ if $a$ is even and the
  opposite one if $a$ is odd. The orientation of $\mathcal{C}$ is
  \emph{coherent} if every ridge lies in exactly two facets and these induce
  opposite orientations on it. A point
  $p = (p_{1}, \ldots, p_{n}) \in (\mathbb{R}^{d})^{n}$ \emph{realises}
  $\mathcal{C}$ if $P = \conv\{p_{1}, \ldots, p_{n}\}$ has the $p_{i}$ as its
  vertices and its facets are exactly the simplices
  $\conv\{p_{j} : j \in F\}$, $F \in \mathcal{F}$. The \emph{signed volume}
  of $p$ is
  \begin{equation}
    \label{eq:signed-volume}
    D(p) = \sum_{F \in \mathcal{F}} \det(p_{F_{1}}, \ldots, p_{F_{d}}).
  \end{equation}
\end{definition}

\begin{lemma}[{\cite[Lemma~3.2]{medgyes-maxvol-d4}}]
  \label{lem:volume}
  If the orientation of $\mathcal{C}$ is coherent and $p$ realises
  $\mathcal{C}$, then
  \[
    d!\,\vol_{d}(\conv\{p_{1}, \ldots, p_{n}\}) = |D(p)|.
  \]
\end{lemma}

\begin{definition}[{\cite[Definition~3.4]{medgyes-maxvol-d4}}]
  \label{def:certificate}
  Let $x_{i} = (x_{i,1}, \ldots, x_{i,d})$ be the coordinates of $p_{i}$, and
  $v(x)$ the vector of all monomials of degree at most $2$ in the $x_{i,a}$. A
  \emph{certificate} for $\mathcal{C}$ and $\beta \in \mathbb{Q}$ is a
  symmetric positive definite matrix $Q$ with rational entries together with
  polynomials $h_{1}, \ldots, h_{n}$ with rational coefficients such that
  \begin{equation}
    \label{eq:certificate}
    \beta - D(x) = v(x)^{\mathsf{T}} Q\, v(x)
      + \sum_{i=1}^{n} h_{i}(x)\bigl(|x_{i}|^{2} - 1\bigr)
  \end{equation}
  holds identically as polynomials in the $x_{i,a}$.
\end{definition}

\begin{proposition}[{\cite[Proposition~3.5]{medgyes-maxvol-d4}}]
  \label{prop:certificate}
  Suppose the orientation of $\mathcal{C}$ is coherent and a certificate for
  $\mathcal{C}$ and $\beta$ exists. Then every polytope inscribed in
  $\mathbb{S}^{d-1}$ that realises $\mathcal{C}$ has volume at most
  $\beta/d!$.
\end{proposition}

A certificate is finite rational data, and its hypotheses (the identity
\eqref{eq:certificate}, positive definiteness of $Q$, coherence, and the facet
sets of the type) are checked in exact arithmetic as described in
\cite[Section~3.3]{medgyes-maxvol-d4}.

\subsection{Simplicial \texorpdfstring{$4$}{4}-polytopes with nine vertices}

\begin{theorem}[{\cite{altshuler-bokowski-steinberg-1980},
    \cite[Table~3]{firsching-2017}}]
  \label{thm:count}
  There are exactly $1142$ combinatorial types of simplicial $4$-polytopes with
  $9$ vertices.
\end{theorem}

The count is due to Altshuler, Bokowski and Steinberg; table numbers of
\cite{firsching-2017} refer to its arXiv version.
We use the realisations of all $1142$ types given there (Lemma~\ref{lem:data}).

\subsection{Direct sums}

\begin{cited}{Direct sum}{{\cite[Section~15.1.3]{henk-richter-gebert-ziegler-2017}}}
  Let $P \subseteq \mathbb{R}^{d}$ be a $d$-polytope with $n$ vertices and $m$
  facets, and $P' \subseteq \mathbb{R}^{d'}$ a $d'$-polytope with $n'$
  vertices and $m'$ facets, both with the origin in their interiors. The
  \emph{direct sum}
  $P \oplus P' = \conv(\{(p, 0) : p \in P\} \cup \{(0, p') : p' \in P'\})$ is
  $(d + d')$-dimensional, and has $n + n'$ vertices and $m \cdot m'$ facets.
\end{cited}

We use this for two regular polygons in orthogonal $2$-planes through the
origin, in Lemma~\ref{lem:freesum}. A polygon has as many edges (its facets)
as vertices, so for polygons with $m$ and $l$ vertices the direct sum has
$m + l$ vertices and $ml$ facets. For $m \geq 3$, $R_{m}$ denotes a regular
$m$-gon inscribed in a unit circle; $\area(R_{m}) = \frac{m}{2}\sin\frac{2\pi}{m}$.

\subsection{Bivectors}

Let $e_{1}, \ldots, e_{N}$ be the standard basis of $\mathbb{R}^{N}$ and
$v \cdot w$ the standard inner product. We use the usual identification of
bivectors with skew-symmetric matrices: the space $\Lambda^{2}\mathbb{R}^{N}$
of \emph{bivectors} is the space of real skew-symmetric $N \times N$ matrices,
and for $v, w \in \mathbb{R}^{N}$ the \emph{wedge product} is
\[
  v \wedge w = vw^{\mathsf{T}} - wv^{\mathsf{T}}.
\]
It is bilinear and satisfies $w \wedge v = -\,v \wedge w$, so $v \wedge v = 0$.
A skew-symmetric matrix $X$ is determined by its entries $X_{kl}$, $k < l$,
above the diagonal, and $e_{kl} = e_{k} \wedge e_{l}$ has entry $1$ at
$(k, l)$, $-1$ at $(l, k)$ and $0$ elsewhere. Hence the $e_{kl}$, $k < l$, form
a basis of $\Lambda^{2}\mathbb{R}^{N}$, we write
$X = \sum_{k<l} X_{kl}\, e_{kl}$, and
\[
  (v \wedge w)_{kl} = v_{k} w_{l} - v_{l} w_{k},
\]
which is the signed area of the projection of the parallelogram spanned by
$v$ and $w$ onto the plane of $e_{k}$ and $e_{l}$. Every bivector is a sum of
wedge products, but for $N \geq 4$ not every bivector is a single wedge
product: $v \wedge w$ has rank at most $2$, while $e_{12} + e_{34}$ has
rank $4$.

We use the inner product
$\langle X, Y\rangle = \frac12\operatorname{tr}(X^{\mathsf{T}}Y)
= \sum_{k<l} X_{kl}Y_{kl}$ on $\Lambda^{2}\mathbb{R}^{N}$, and
$|X| = \langle X, X\rangle^{1/2}$; the factor $\frac12$ accounts for each pair
$k < l$ appearing twice in $\operatorname{tr}(X^{\mathsf{T}}Y) = \sum_{i,j}
X_{ij}Y_{ij}$. Expanding the trace with
$\operatorname{tr}\bigl((pq^{\mathsf{T}})^{\mathsf{T}} rs^{\mathsf{T}}\bigr)
= (p \cdot r)(q \cdot s)$ gives
\[
  \langle a \wedge b, c \wedge d\rangle = (a \cdot c)(b \cdot d) - (a \cdot d)(b \cdot c).
\]
By this formula, $u_{k} \wedge u_{l}$, $k < l$, are orthonormal for every
orthonormal basis $u_{1}, \ldots, u_{N}$ of $\mathbb{R}^{N}$. Taking $c = a$
and $d = b$ in the formula gives the following.

\begin{lemma}[Lagrange's identity]
  \label{lem:lagrange}
  For $a, b \in \mathbb{R}^{N}$,
  $|a \wedge b|^{2} = |a|^{2}|b|^{2} - (a \cdot b)^{2}$.
\end{lemma}

We use Lemma~\ref{lem:lagrange} in the form $|a \wedge b| \leq |a|\,|b| \leq
\frac12(|a|^{2} + |b|^{2})$, where both inequalities are equalities if and only
if $a \cdot b = 0$ and $|a| = |b|$.

Let $\Pi \subseteq \mathbb{R}^{N}$ be a $2$-plane through the origin and
$\Pi^{\perp}$ its orthogonal complement. For an orthonormal basis $u_{1},
u_{2}$ of $\Pi$, the \emph{unit area element} of $\Pi$ is
$\omega_{\Pi} = u_{1} \wedge u_{2}$; for instance, $e_{kl}$ is that of the
plane spanned by $e_{k}$ and $e_{l}$. By bilinearity and antisymmetry,
changing the orthonormal basis multiplies $u_{1} \wedge u_{2}$ by the
determinant of the change of basis, which is $\pm 1$; statements about
$\omega_{\Pi}$ below hold for either sign. By Lemma~\ref{lem:lagrange}, $|\omega_{\Pi}| = 1$.

\subsection{The Hodge star on \texorpdfstring{$\Lambda^{2}\mathbb{R}^{4}$}{Lambda2R4}}

We use the Hodge star only for $N = 4$. For $k < l$ in $\{1, 2, 3, 4\}$ let
$\{k', l'\}$, $k' < l'$, be the complement of $\{k, l\}$, so that the plane
spanned by $e_{k'}, e_{l'}$ is the orthogonal complement of the plane spanned
by $e_{k}, e_{l}$. Let $\varepsilon_{kl} = \pm 1$ be the sign of the
permutation $(k, l, k', l')$ of $(1, 2, 3, 4)$. The \emph{Hodge star} is the
linear map $\star$ of $\Lambda^{2}\mathbb{R}^{4}$ with
\[
  \star e_{kl} = \varepsilon_{kl}\, e_{k'l'},
\]
that is, it sends the unit area element of each coordinate plane to that of its
orthogonal complement, up to sign:
\[
  \star e_{12} = e_{34}, \quad \star e_{34} = e_{12}, \quad
  \star e_{14} = e_{23}, \quad \star e_{23} = e_{14}, \quad
  \star e_{13} = -e_{24}, \quad \star e_{24} = -e_{13}.
\]
It permutes the orthonormal basis $e_{kl}$ up to signs, so it is an isometry.
Its use here is that, by Laplace expansion, it turns a $4 \times 4$
determinant into an inner product of bivectors.

\subsection{The Laplace expansion}

\begin{theorem}[{The Laplace expansion theorem, \cite[Section~0.8.9]{horn-johnson-2013}}]
  \label{thm:laplace}
  Let $A \in M_{n}(F)$, let $k \in \{1, \ldots, n\}$ be given, and let
  $\beta \subseteq \{1, \ldots, n\}$ be any given index set of cardinality $k$.
  Then
  \[
    \det A = \sum_{\alpha} (-1)^{p(\alpha,\beta)} \det A[\alpha, \beta]
      \det A[\alpha^{c}, \beta^{c}]
    = \sum_{\alpha} (-1)^{p(\alpha,\beta)} \det A[\beta, \alpha]
      \det A[\beta^{c}, \alpha^{c}]
  \]
  in which the sums are over all index sets $\alpha \subseteq \{1, \ldots, n\}$
  of cardinality $k$, and
  $p(\alpha, \beta) = \sum_{i \in \alpha} i + \sum_{j \in \beta} j$.
\end{theorem}

Here $M_{n}(F)$ denotes the $n \times n$ matrices over a field $F$,
$A[\alpha, \beta]$ is the submatrix of $A$ with rows indexed by $\alpha$ and
columns indexed by $\beta$, and $\alpha^{c}$, $\beta^{c}$ are the complements
of $\alpha$, $\beta$ \cite[Section~0.7.1]{horn-johnson-2013}. We use the first
expansion of Theorem~\ref{thm:laplace} for the $4 \times 4$ matrix $A$ with columns $v, v', w, w'$, index
sets of cardinality $2$, and $\beta = \{1, 2\}$, the columns of $v, v'$. For
$\alpha = \{k, l\}$, $k < l$, with complement $\{k', l'\}$, $k' < l'$, the
minor $\det A[\alpha, \beta]$ is $v_{k}v'_{l} - v_{l}v'_{k} =
(v \wedge v')_{kl}$, the complementary minor $\det A[\alpha^{c}, \beta^{c}]$
is $(w \wedge w')_{k'l'}$, and $(-1)^{p(\alpha, \beta)} = (-1)^{k + l + 1 + 2}$,
which is the sign $\varepsilon_{kl}$ of the permutation $(k, l, k', l')$.
Hence
\[
  \det(v, v', w, w')
  = \sum_{k<l} \varepsilon_{kl}\,(v \wedge v')_{kl}\,(w \wedge w')_{k'l'}.
\]
Since exchanging the two pairs is an even permutation,
$\varepsilon_{k'l'} = \varepsilon_{kl}$, so the coefficient of $e_{kl}$ in
$\star Y$ is $\varepsilon_{kl} Y_{k'l'}$, and this sum is an inner product:
\begin{equation}
  \label{eq:detstar}
  \det(v, v', w, w') = \langle v \wedge v', \star(w \wedge w')\rangle.
\end{equation}
For example, $\det(e_{1}, e_{2}, e_{3}, e_{4}) = \langle e_{12}, \star
e_{34}\rangle = \langle e_{12}, e_{12}\rangle = 1$. So, up to the areas
$|v \wedge v'|$ and $|w \wedge w'|$, the determinant compares the plane of
$v, v'$ with the orthogonal complement of that of $w, w'$.

\begin{remark}
  The description of $\star$ on coordinate planes holds for every plane: for
  a $2$-plane $\Pi \subseteq \mathbb{R}^{4}$ through the origin,
  \begin{equation}
    \label{eq:starplane}
    \star\,\omega_{\Pi} = \pm\,\omega_{\Pi^{\perp}}.
  \end{equation}
  Indeed, extend an orthonormal basis $u_{1}, u_{2}$ of $\Pi$ by one $u_{3},
  u_{4}$ of $\Pi^{\perp}$. By \eqref{eq:detstar}, $\langle u_{i} \wedge u_{j},
  \star(u_{1} \wedge u_{2})\rangle = \det(u_{i}, u_{j}, u_{1}, u_{2})$, which
  vanishes unless $\{i, j\} = \{3, 4\}$ and is $\pm 1$ then. The six
  $u_{i} \wedge u_{j}$, $i < j$, are orthonormal, hence a basis of
  $\Lambda^{2}\mathbb{R}^{4}$, and these inner products are the coordinates
  of $\star\omega_{\Pi}$ in it; so $\star\omega_{\Pi} = \pm\, u_{3} \wedge
  u_{4} = \pm\,\omega_{\Pi^{\perp}}$.
\end{remark}

\section{Lemmas on the join of two cycles}
\label{sec:lemmas}

For $m \geq 3$ and $x = (x_{1}, \ldots, x_{m}) \in (\mathbb{R}^{N})^{m}$, with
indices modulo $m$, let
\begin{equation}
  \label{eq:areavector}
  \mathcal{A}(x) = \sum_{j=1}^{m} x_{j} \wedge x_{j+1}.
\end{equation}
Thus $\frac12\mathcal{A}(x)$ is the \emph{area vector} of the closed polygon
$x_{1} \cdots x_{m}$; for $N = 2$ its single coordinate is the signed area given
by the shoelace formula.

\begin{definition}
  \label{def:join}
  For $m, l \geq 3$ let $\mathcal{J}_{m,l} = (m + l, \mathcal{F}_{m,l})$ with
  \[
    \mathcal{F}_{m,l} = \bigl\{(i,\ i+1,\ m+j,\ m+j+1) : i \in \mathbb{Z}_{m},\
      j \in \mathbb{Z}_{l}\bigr\},
  \]
  where the first two entries are read modulo $m$ in $\{1, \ldots, m\}$ and the
  last two modulo $l$ in $\{m+1, \ldots, m+l\}$. For
  $p \in (\mathbb{R}^{4})^{m+l}$ write $a_{i} = p_{i}$, $b_{j} = p_{m+j}$, and
  \[
    A = \mathcal{A}(a_{1}, \ldots, a_{m}), \qquad
    B = \mathcal{A}(b_{1}, \ldots, b_{l}).
  \]
\end{definition}

Each facet of $\mathcal{J}_{m,l}$ consists of an edge of the cycle $(1, \ldots,
m)$ and an edge of the cycle $(m+1, \ldots, m+l)$, so as a simplicial complex
$\mathcal{J}_{m,l}$ is the join of these two cycles. By
Lemma~\ref{lem:freesum} it is the boundary complex of the direct sum of two
regular polygons; the polytope itself is this direct sum, not a join of
polytopes.

By \eqref{eq:signed-volume}
\begin{equation}
  \label{eq:Djoin}
  D(p) = \sum_{i \in \mathbb{Z}_{m}}\sum_{j \in \mathbb{Z}_{l}}
    \det(a_{i}, a_{i+1}, b_{j}, b_{j+1}).
\end{equation}

\begin{lemma}
  \label{lem:freesum}
  Let $m, l \geq 3$, let $u_{1}, u_{2}, u_{3}, u_{4}$ be an orthonormal basis of
  $\mathbb{R}^{4}$, and for $i \in \mathbb{Z}_{m}$, $j \in \mathbb{Z}_{l}$ let
  \[
    a_{i} = \cos\tfrac{2\pi i}{m}\,u_{1} + \sin\tfrac{2\pi i}{m}\,u_{2}, \qquad
    b_{j} = \cos\tfrac{2\pi j}{l}\,u_{3} + \sin\tfrac{2\pi j}{l}\,u_{4}.
  \]
  Then $P = \conv\{a_{i}, b_{j}\} \in \mathcal{P}_{4}(m + l)$, and the point
  $p = (a_{1}, \ldots, a_{m}, b_{1}, \ldots, b_{l})$ realises
  $\mathcal{J}_{m,l}$.
\end{lemma}

\begin{proof}
  In the coordinates of the basis $u_{1}, \ldots, u_{4}$, $P$ is the direct sum
  of two regular polygons with $m$ and $l$ vertices and edges, inscribed in unit
  circles centred at the origin. By the cited statement on direct sums, $P$ is
  $4$-dimensional, all $m + l$ points $a_{i}, b_{j}$ are vertices, and $P$ has
  $ml$ facets. All the points are unit vectors, so $P \in \mathcal{P}_{4}(m +
  l)$.

  It remains to show that each of the $ml$ distinct simplices
  $\conv\{a_{i}, a_{i+1}, b_{j}, b_{j+1}\}$ is a facet. Let $h_{m} =
  \cos\frac{\pi}{m}$ and $\nu_{i} = \cos\frac{(2i+1)\pi}{m}\,u_{1} +
  \sin\frac{(2i+1)\pi}{m}\,u_{2}$, so that $\nu_{i} \cdot a_{k} =
  \cos\frac{(2k-2i-1)\pi}{m} \leq h_{m}$, with equality exactly for $k \in
  \{i, i+1\}$, and $\nu_{i} \cdot b_{k} = 0$; define $h_{l}$ and $\mu_{j}$ in the
  same way for the $b$'s. The linear function $\varphi(x) = \nu_{i} \cdot
  x/h_{m} + \mu_{j} \cdot x/h_{l}$ is at most $1$ at every vertex of $P$, with
  equality exactly at $a_{i}, a_{i+1}, b_{j}, b_{j+1}$. So
  $P \cap \{\varphi = 1\} = \conv\{a_{i}, a_{i+1}, b_{j}, b_{j+1}\}$ is a face,
  and it is $3$-dimensional, since the four points are linearly independent.
  Since $P$ has exactly $ml$ facets, these are all of them, so $p$ realises
  $\mathcal{J}_{m,l}$.
\end{proof}

\begin{lemma}
  \label{lem:coherent}
  The orientation of $\mathcal{J}_{m,l}$ is coherent.
\end{lemma}

\begin{proof}
  Every ridge is of one of two kinds. The ridge $\{a_{i}, b_{j}, b_{j+1}\}$
  lies in exactly the facets $(a_{i-1}, a_{i}, b_{j}, b_{j+1})$ and $(a_{i},
  a_{i+1}, b_{j}, b_{j+1})$, which differ as $m \geq 3$. Deleting entry~$1$ of
  the first gives the opposite of $(a_{i}, b_{j}, b_{j+1})$, and deleting
  entry~$2$ of the second gives $(a_{i}, b_{j}, b_{j+1})$. The ridge $\{a_{i},
  a_{i+1}, b_{j}\}$ lies in exactly $(a_{i}, a_{i+1}, b_{j-1}, b_{j})$ and
  $(a_{i}, a_{i+1}, b_{j}, b_{j+1})$, and deleting entry~$3$, respectively
  entry~$4$, gives opposite orientations in the same way.
\end{proof}

\begin{lemma}
  \label{lem:factor}
  For every $p \in (\mathbb{R}^{4})^{m+l}$, with $D(p)$ as in \eqref{eq:Djoin}
  and $A$, $B$ as in Definition~\ref{def:join},
  $D(p) = \langle A, \star B\rangle$. Hence $|D(p)| \leq |A|\,|B|$, and if
  $A, B \neq 0$, equality holds if and only if $\star B = \lambda A$ for some
  $\lambda \in \mathbb{R}$.
\end{lemma}

\begin{proof}
  By \eqref{eq:detstar}, each term of \eqref{eq:Djoin} is
  $\langle a_{i} \wedge a_{i+1}, \star(b_{j} \wedge b_{j+1})\rangle$; summing
  over $i$ and $j$ gives $D(p) = \langle A, \star B\rangle$. The rest is the
  Cauchy--Schwarz inequality, with $|\star B| = |B|$.
\end{proof}

\section{An inequality for closed polygons}
\label{sec:polygon}

By Lemma~\ref{lem:factor}, $|D(p)| \leq |A|\,|B|$, where $A$ and $B$ are the
values of $\mathcal{A}$, defined in \eqref{eq:areavector}, on the two cycles of
Definition~\ref{def:join}. So a sharp bound on $D$
reduces to a sharp bound on $|\mathcal{A}(x)|$ for a single cycle of unit
vertices, together with its equality case. Such a polygon in $\mathbb{R}^{4}$
is in general neither planar nor convex, so the classical bound for polygons
inscribed in a circle does not apply. Proposition~\ref{prop:polygon} gives the
bound for $m \leq 5$, the cycle lengths that occur.

\begin{notation}
  For $m \geq 3$ and $x = (x_{1}, \ldots, x_{m}) \in (\mathbb{R}^{N})^{m}$, with
  indices modulo $m$ and $\mathcal{A}(x)$ as in \eqref{eq:areavector}, let
  $\theta = 2\pi/m$ and, for $1 \leq k < m/2$,
  \[
    y_{k} = \sqrt{2/m}\,\sum_{j=1}^{m} \cos(jk\theta)\, x_{j}, \qquad
    z_{k} = \sqrt{2/m}\,\sum_{j=1}^{m} \sin(jk\theta)\, x_{j}.
  \]
\end{notation}

The vectors $y_{k}, z_{k} \in \mathbb{R}^{N}$ are the
coefficients of $x$ in the real discrete Fourier basis, and $\sin k\theta > 0$
for $1 \leq k < m/2$.

Since $\mathcal{A}(x)$ is quadratic in $x$ and unchanged by a cyclic shift of
the indices, the discrete Fourier transform diagonalises it.
Lemma~\ref{lem:fourier} makes this explicit: it writes $\mathcal{A}(x)$ as a
weighted sum of the bivectors $y_{k} \wedge z_{k}$, one for each frequency $k$,
and bounds the total size of the coefficients by $\sum_{j}|x_{j}|^{2}$, which
is $m$ for unit vertices. For $N = 2$, half of \eqref{eq:fourier} is, in real
form, Schoenberg's expression for the oriented area of a plane polygon through
its finite Fourier coefficients \cite[Section~8, (38)]{schoenberg-1950}.
Proposition~\ref{prop:polygon} then bounds each term,
and the weight $\sin k\theta$ is uniquely largest at $k = 1$, the frequency of
the regular $m$-gon, exactly when $m \leq 5$ (Remark~\ref{rem:polygon}).

\begin{lemma}
  \label{lem:fourier}
  Let $m \geq 3$ and $x \in (\mathbb{R}^{N})^{m}$. Then
  \begin{equation}
    \label{eq:fourier}
    \mathcal{A}(x) = \sum_{1 \leq k < m/2} 2\sin k\theta\; y_{k} \wedge z_{k}
  \end{equation}
  and
  \begin{equation}
    \label{eq:parseval}
    \sum_{1 \leq k < m/2}\bigl(|y_{k}|^{2} + |z_{k}|^{2}\bigr)
      \leq \sum_{j=1}^{m}|x_{j}|^{2},
  \end{equation}
  with equality if and only if
  \begin{equation}
    \label{eq:inversion}
    x_{j} = \sqrt{2/m}\,\sum_{1 \leq k < m/2}\bigl(\cos(jk\theta)\, y_{k}
      + \sin(jk\theta)\, z_{k}\bigr) \qquad (j = 1, \ldots, m).
  \end{equation}
\end{lemma}

\begin{proof}
  Let $f_{k}, g_{k} \in \mathbb{R}^{m}$ have entries $f_{k,j} =
  \sqrt{2/m}\,\cos jk\theta$ and $g_{k,j} = \sqrt{2/m}\,\sin jk\theta$ for
  $1 \leq k < m/2$, and $f_{k,j} = \cos(jk\theta)/\sqrt{m}$ for $k = 0$ and, if
  $m$ is even, $k = m/2$. These $m$ vectors form an orthonormal basis
  $\mathcal{B}$ of $\mathbb{R}^{m}$: by the product formulas for $\cos$ and
  $\sin$, the inner product of two of them is a combination of the sums
  $\sum_{j} \cos jn\theta$ and $\sum_{j} \sin jn\theta$. Both vanish for
  $n \not\equiv 0 \pmod{m}$ \cite[Chapter~7, Section~1]{stein-shakarchi-2003},
  the sine sum vanishes termwise for $n \equiv 0$, and a cosine sum with
  $n \equiv 0$ occurs only in the squared lengths.

  For $v \in \mathcal{B}$ let $x_{v} = \sum_{j} v_{j}\, x_{j} \in
  \mathbb{R}^{N}$, so that $x_{f_{k}} = y_{k}$ and $x_{g_{k}} = z_{k}$.
  Expanding each coordinate of $x$ in $\mathcal{B}$ gives $x_{j} =
  \sum_{v \in \mathcal{B}} v_{j}\, x_{v}$, and Parseval's identity gives
  $\sum_{j}|x_{j}|^{2} = \sum_{v \in \mathcal{B}}|x_{v}|^{2}$. This proves
  \eqref{eq:parseval}, with equality if and only if $x_{f_{0}} = 0$ and, for
  even $m$, $x_{f_{m/2}} = 0$, that is, if and only if \eqref{eq:inversion}
  holds.

  Let $K$ be the $m \times m$ matrix with $(Kv)_{j} = v_{j+1} - v_{j-1}$.
  Substituting the expansion of $x_{j}$ and $x_{j+1}$ into $\mathcal{A}(x)$,
  and using $x_{v'} \wedge x_{v} = -\,x_{v} \wedge x_{v'}$ and a shift of the
  index $j$,
  \[
    \mathcal{A}(x) = \sum_{v, v' \in \mathcal{B}}\Bigl(\sum_{j} v_{j}
      v'_{j+1}\Bigr)\, x_{v} \wedge x_{v'}
    = \frac12\sum_{v, v' \in \mathcal{B}} \bigl(v^{\mathsf{T}} K v'\bigr)\,
      x_{v} \wedge x_{v'}.
  \]
  Since $\cos(j+1)\varphi - \cos(j-1)\varphi = -2\sin\varphi\,\sin j\varphi$
  and $\sin(j+1)\varphi - \sin(j-1)\varphi = 2\sin\varphi\,\cos j\varphi$, we
  have $Kf_{k} = -2\sin k\theta\; g_{k}$ and $Kg_{k} = 2\sin k\theta\; f_{k}$
  for $1 \leq k < m/2$, while $K$ vanishes on $f_{0}$ and $f_{m/2}$. So
  $v^{\mathsf{T}} K v'$ is nonzero only for $(v, v') = (f_{k}, g_{k})$, where it
  is $2\sin k\theta$, and for $(v, v') = (g_{k}, f_{k})$, where it is
  $-2\sin k\theta$. Together these two terms give $2\sin k\theta\; y_{k} \wedge
  z_{k}$, which proves \eqref{eq:fourier}.
\end{proof}

\begin{proposition}
  \label{prop:polygon}
  Let $m \in \{3, 4, 5\}$ and $x \in (\mathbb{R}^{N})^{m}$ with $|x_{j}| = 1$
  for all $j$. Then $|\mathcal{A}(x)| \leq m\sin\frac{2\pi}{m}$, with equality if
  and only if
  \begin{equation}
    \label{eq:regular}
    x_{j} = \cos\frac{2\pi j}{m}\,u + \sin\frac{2\pi j}{m}\,w
      \qquad (j = 1, \ldots, m)
  \end{equation}
  for orthonormal $u, w$, that is, if the $x_{j}$ are the consecutive vertices
  of a regular $m$-gon inscribed in the unit circle of $\Pi = \spn(u, w)$. In
  this case $\mathcal{A}(x) = m\sin\frac{2\pi}{m}\; u \wedge w$.
\end{proposition}

\begin{proof}
  The indices $1 \leq k < m/2$ are $k = 1$ and, for $m = 5$, $k = 2$, with
  $\sin 2\theta = \sin\frac{4\pi}{5} < \sin\frac{2\pi}{5} = \sin\theta$. By
  \eqref{eq:fourier}, the triangle inequality, Lemma~\ref{lem:lagrange} in the
  form $|a \wedge b| \leq \frac12(|a|^{2} + |b|^{2})$, $\sin k\theta \leq
  \sin\theta$, and \eqref{eq:parseval},
  \[
    |\mathcal{A}(x)|
      \leq \sum_{k} 2\sin k\theta\; |y_{k} \wedge z_{k}|
      \leq \sum_{k} \sin k\theta\,\bigl(|y_{k}|^{2} + |z_{k}|^{2}\bigr)
      \leq \sin\theta\,\sum_{k}\bigl(|y_{k}|^{2} + |z_{k}|^{2}\bigr)
      \leq \sin\theta\,\sum_{j}|x_{j}|^{2} = m\sin\theta.
  \]

  Suppose equality holds, so that each inequality above is an equality. By the
  last one and Lemma~\ref{lem:fourier}, \eqref{eq:inversion} holds and
  $\sum_{k}(|y_{k}|^{2} + |z_{k}|^{2}) = m$. For $m = 5$, the third gives
  $(\sin\theta - \sin 2\theta)(|y_{2}|^{2} + |z_{2}|^{2}) = 0$, so $y_{2} =
  z_{2} = 0$. Hence, for every $m \in \{3, 4, 5\}$, $|y_{1}|^{2} + |z_{1}|^{2}
  = m$, and \eqref{eq:inversion} reads $x_{j} = \sqrt{2/m}\,(\cos(j\theta)\,
  y_{1} + \sin(j\theta)\, z_{1})$. By the second inequality and the equality
  condition stated after Lemma~\ref{lem:lagrange}, $y_{1} \cdot z_{1} = 0$ and $|y_{1}| =
  |z_{1}| = \sqrt{m/2}$. Thus $u = \sqrt{2/m}\,y_{1}$ and $w =
  \sqrt{2/m}\,z_{1}$ are orthonormal and \eqref{eq:regular} holds.

  Conversely, if \eqref{eq:regular} holds, then $x_{j} \wedge x_{j+1} = (\cos
  j\theta \sin(j+1)\theta - \sin j\theta \cos(j+1)\theta)\, u \wedge w =
  \sin\theta\; u \wedge w$ for every $j$, so $\mathcal{A}(x) = m\sin\theta\;
  u \wedge w$, which has length $m\sin\theta$ by Lemma~\ref{lem:lagrange}.
\end{proof}

\begin{remark}
  \label{rem:polygon}
  For $N = 2$, Proposition~\ref{prop:polygon} is the case $m \leq 5$ of
  Schoenberg's bound for plane polygons
  \cite[(39) and Theorem~5]{schoenberg-1950}, applied to the polygon and to its
  reverse. Schoenberg bounds the oriented area, which is
  $\frac12\mathcal{A}(x)$, by $\frac12\max_{k}\sin k\theta\,
  \sum_{j}|x_{j}|^{2}$; in our normalisation the bound reads
  $\mathcal{A}(x) \leq \max_{k}\sin k\theta\, \sum_{j}|x_{j}|^{2}$. For unit
  vectors in counterclockwise order it is the classical fact that
  among the $m$-gons inscribed in a circle the regular one has the greatest
  area \cite[Theorem~12]{kazarinoff-1961}. For $N = 3$ the bound also follows
  from the planar case, since every bivector in $\Lambda^{2}\mathbb{R}^{3}$ is a
  multiple of some $\omega_{\Pi}$, and $\mathcal{A}(\pi x) = \langle
  \mathcal{A}(x), \omega_{\Pi}\rangle\,\omega_{\Pi}$ for the orthogonal
  projection $\pi$ onto $\Pi$, which does not increase
  $\sum_{j}|x_{j}|^{2}$. For $N \geq 4$ this reduction fails: projecting onto planes
  bounds only $\max_{\Pi}\langle \mathcal{A}(x), \omega_{\Pi}\rangle$, which
  can be smaller than $|\mathcal{A}(x)|$, as for $e_{12} + e_{34}$, where it
  is $1$ while the length is $\sqrt{2}$. Lemma~\ref{lem:factor} needs the
  length itself. The proof uses $|x_{j}| = 1$ only through
  $\sum_{j}|x_{j}|^{2} = m$.

  The restriction on $m$ is needed. For every $m \geq 3$ the chain above, with
  $\max_{k}\sin k\theta$ in place of $\sin\theta$, gives $|\mathcal{A}(x)| \leq
  m\max_{k}\sin k\theta$, and the computation in the converse, with $k\theta$ in
  place of $\theta$, shows that $x_{j} = \cos(jk\theta)\,u + \sin(jk\theta)\,w$
  attains $m\sin k\theta$. For $k = 2$ these points traverse a regular star
  polygon if $m$ is odd and a regular $\frac{m}{2}$-gon twice if $m$ is even.
  For $m = 6$ this ties with the hexagon, and for $m \geq 7$ it beats the
  regular $m$-gon, since $\sin 2\theta > \sin\theta$ when $\theta < \pi/3$.
  For $N = 2$, Schoenberg determines the extremal polygons for every $m$
  \cite[Theorem~5]{schoenberg-1950}.
\end{remark}

\section{The optimum within the join type}
\label{sec:pair}

\begin{theorem}
  \label{thm:pair}
  Let $m, l \in \{3, 4, 5\}$, $p \in (\mathbb{S}^{3})^{m+l}$, and $D(p)$ the
  signed volume \eqref{eq:Djoin} of $\mathcal{J}_{m,l}$, with $a_{i}$, $b_{j}$
  as in Definition~\ref{def:join}. Then
  \begin{equation}
    \label{eq:pairbound}
    |D(p)| \leq m\sin\frac{2\pi}{m}\cdot l\sin\frac{2\pi}{l},
  \end{equation}
  with equality if and only if, for some $2$-plane $\Pi$, the points $a_{1},
  \ldots, a_{m}$ are the consecutive vertices of a regular $m$-gon inscribed in
  the unit circle of $\Pi$ and $b_{1}, \ldots, b_{l}$ those of a regular
  $l$-gon inscribed in the unit circle of $\Pi^{\perp}$.
\end{theorem}

\begin{proof}
  With $A$ and $B$ as in Definition~\ref{def:join},
  Lemma~\ref{lem:factor} gives $|D(p)| \leq |A|\,|B|$, and by
  Proposition~\ref{prop:polygon}, $|A| \leq m\sin\frac{2\pi}{m}$ and $|B| \leq
  l\sin\frac{2\pi}{l}$. This proves \eqref{eq:pairbound}.

  If equality holds, then $A, B \neq 0$ and equality holds in all three
  inequalities. By Proposition~\ref{prop:polygon}, the $a_{i}$ are as in
  \eqref{eq:regular} in a plane $\Pi$, with $A = m\sin\frac{2\pi}{m}\,
  \omega_{\Pi}$, and the $b_{j}$ in a plane $\Pi'$, with $B =
  l\sin\frac{2\pi}{l}\,\omega_{\Pi'}$. By Lemma~\ref{lem:factor}, $\star B =
  \lambda A$, so $\star\omega_{\Pi'}$ is a multiple of $\omega_{\Pi}$; both
  have length $1$, as $\star$ is an isometry, so by \eqref{eq:starplane}
  $\omega_{\Pi'^{\perp}} = \pm\star\omega_{\Pi'} = \pm\omega_{\Pi}$. With
  orthonormal bases $u_{1}, u_{2}$ of $\Pi$ and $w_{1},
  w_{2}$ of $\Pi'^{\perp}$, $\omega_{\Pi} = u_{1}u_{2}^{\mathsf{T}} -
  u_{2}u_{1}^{\mathsf{T}}$ has column space $\Pi$, and likewise for
  $\Pi'^{\perp}$; hence $\Pi = \Pi'^{\perp}$.

  Conversely, let $u_{1}, u_{2}$ and $u_{3}, u_{4}$ be orthonormal bases of
  $\Pi$ and $\Pi^{\perp}$ for which the $a_{i}$ and $b_{j}$ have the form
  \eqref{eq:regular}. By Proposition~\ref{prop:polygon}, $A = m\sin\frac{2\pi}{m}\,u_{1}
  \wedge u_{2}$ and $B = l\sin\frac{2\pi}{l}\,u_{3} \wedge u_{4}$, and by
  Lemma~\ref{lem:factor} and \eqref{eq:starplane}, $|D(p)| = |\langle A, \star
  B\rangle| = |A|\,|B|$.
\end{proof}

\begin{remark}
  \label{rem:bipyramid}
  The boundary complex of a double $n$-pyramid in $\mathbb{R}^{3}$ is the join
  of an $(n-2)$-cycle and a pair of points, and Berman and Hanes showed that an inscribed double
  $n$-pyramid with Property Z is unique up to congruence, with volume
  $\frac{n-2}{3}\sin\frac{2\pi}{n-2}$ \cite[Lemma~2]{berman-hanes-1970}; for
  $d$-bipyramids, Horv\'ath and L\'angi showed that Property Z forces the
  apexes to be antipodal and the base to lie in the orthogonal hyperplane
  \cite[Corollary~2]{horvath-langi-2016}. Theorem~\ref{thm:pair} is the
  analogue in $\mathbb{R}^{4}$ for $\mathcal{J}_{m,l}$, the join of two
  cycles. Unlike these
  results, which use Property Z, it bounds $|D(p)|$ for every
  $p \in (\mathbb{S}^{3})^{m+l}$.
\end{remark}

\begin{corollary}
  \label{cor:pair}
  Let $m, l \in \{3, 4, 5\}$. If $P \in \mathcal{P}_{4}(m+l)$ is
  combinatorially equivalent to the polytope of Lemma~\ref{lem:freesum}, then
  \[
    \vol_{4}(P) \leq \frac{1}{24}\, m\sin\frac{2\pi}{m}\; l\sin\frac{2\pi}{l}
      = \frac{1}{6}\area(R_{m})\area(R_{l}),
  \]
  with equality if and only if $P$ is the polytope of Lemma~\ref{lem:freesum}
  for some orthonormal basis $u_{1}, \ldots, u_{4}$, that is, up to an
  orthogonal transformation, the convex hull of $R_{m}$ and $R_{l}$ inscribed
  in $\mathbb{S}^{3}$ in orthogonal $2$-planes through the origin.
\end{corollary}

\begin{proof}
  Label the vertices of $P$ so that $p \in (\mathbb{S}^{3})^{m+l}$ realises
  $\mathcal{J}_{m,l}$. By Lemmas~\ref{lem:coherent} and~\ref{lem:volume},
  $24\vol_{4}(P) = |D(p)|$, and Theorem~\ref{thm:pair} gives the bound. In the
  case of equality, Theorem~\ref{thm:pair} and \eqref{eq:regular} give the
  basis $u_{1}, \ldots, u_{4}$. Conversely, by Lemma~\ref{lem:freesum} the
  polytope of that lemma lies in $\mathcal{P}_{4}(m+l)$ and its vertices, in
  the order given there, realise $\mathcal{J}_{m,l}$; they have the form of
  the equality case of Theorem~\ref{thm:pair}, so its volume is
  $\frac{1}{24}\, m\sin\frac{2\pi}{m}\; l\sin\frac{2\pi}{l}$.
\end{proof}

\section{The other combinatorial types}
\label{sec:rivals}

Firsching's data \cite{firsching-2017} give a polytope with rational
vertices for each line of a file with $1142$ lines \artref{data}; we label
them \texttt{firsching-9-0} to \texttt{firsching-9-1141}.

\begin{lemma}
  \label{lem:data}
  The polytopes \texttt{firsching-9-0} to \texttt{firsching-9-1141} represent
  every combinatorial type of simplicial $4$-polytope with nine vertices, each
  exactly once, and the type of \texttt{firsching-9-10} is that of the
  polytope of Lemma~\ref{lem:freesum} with $(m, l) = (4, 5)$.
\end{lemma}

\begin{proof}
  We check in exact arithmetic, as in \cite[Remark~2.14]{medgyes-maxvol-d4}
  for eight vertices, that these $1142$ polytopes are simplicial and pairwise
  not combinatorially equivalent \artref{testnine}. By
  Theorem~\ref{thm:count} they represent every type. The facets of
  \texttt{firsching-9-10} are exactly those of $\mathcal{J}_{4,5}$ for the
  $4$-cycle $(1, 3, 6, 4)$ and the $5$-cycle $(0, 2, 8, 7, 5)$ of its
  zero-based vertex labels \artref{testnine}. By Lemma~\ref{lem:freesum}, the
  facets of the polytope there are those of $\mathcal{J}_{4,5}$ too, so the two
  are combinatorially equivalent.
\end{proof}

\begin{proposition}
  \label{prop:rivals}
  Every simplicial $P \in \mathcal{P}_{4}(9)$ that is not combinatorially
  equivalent to \texttt{firsching-9-10} has $\vol_{4}(P) \leq 39/50$.
\end{proposition}

\begin{proof}
  By Lemma~\ref{lem:data}, $P$ is combinatorially equivalent to
  \texttt{firsching-9-$k$} for some $k \neq 10$. For each of these $1141$
  values of $k$, the stored
  certificate with $\beta = 468/25$ passes the checks of
  \cite[Section~3.3]{medgyes-maxvol-d4}, including that its facets are those
  of \texttt{firsching-9-$k$} \artref{certificates}, \artref{verifier},
  \artref{testnine}. After
  relabelling its vertices, $P$ realises the coherently oriented complex of
  that certificate, and Proposition~\ref{prop:certificate} gives
  $\vol_{4}(P) \leq \beta/4! = 39/50$.
\end{proof}

\section{The main theorem}
\label{sec:main}

\begin{theorem}
  \label{thm:main}
  \[
    v_{4}(9) = \frac{5}{6}\sin\frac{2\pi}{5}
      = \frac{5}{24}\sqrt{10 + 2\sqrt{5}} = 0.792547\ldots,
  \]
  and every $P \in \mathcal{P}_{4}(9)$ with $\vol_{4}(P) = v_{4}(9)$ is, up to
  an orthogonal transformation, the convex hull of a square and a regular
  pentagon inscribed in $\mathbb{S}^{3}$ in orthogonal $2$-planes through the
  origin.
\end{theorem}

\begin{proof}
  Let $V^{*} = \frac{1}{24}\cdot 4\sin\frac{2\pi}{4}\cdot 5\sin\frac{2\pi}{5} =
  \frac{5}{6}\sin\frac{2\pi}{5}$; since $\sin\frac{2\pi}{5} =
  \frac14\sqrt{10 + 2\sqrt5}$, this is the stated value. By
  Lemma~\ref{lem:freesum} and Corollary~\ref{cor:pair} with $(m, l) = (4, 5)$,
  the square--pentagon pair lies in $\mathcal{P}_{4}(9)$ and has volume
  $V^{*}$, so $v_{4}(9) \geq V^{*}$.

  Let $P \in \mathcal{P}_{4}(9)$ with $\vol_{4}(P) = v_{4}(9)$, which exists by
  Remark~\ref{rem:attained}. By Remark~\ref{rem:propertyz} and
  Lemma~\ref{lem:simplicial}, $P$ is simplicial. Since $39/50 < V^{*}$,
  Proposition~\ref{prop:rivals} shows that $P$ has the combinatorial type of
  \texttt{firsching-9-10}, which by Lemma~\ref{lem:data} is that of the
  polytope of Lemma~\ref{lem:freesum} with $(m, l) = (4, 5)$. Corollary~\ref{cor:pair} gives
  $v_{4}(9) = \vol_{4}(P) \leq V^{*}$, so $v_{4}(9) = V^{*}$, and its equality
  case identifies $P$.
\end{proof}


\section*{AI assistance}

An AI assistant was used during this work, including for mathematical
discussion, drafting and revising the text, and checking the cited statements
against their sources. The proofs do not depend on the assistant: they rest on
the cited results, on the arguments given here, and on the certificates, which
the exact verifier checks.

\section*{Code and data}
\label{sec:artifacts}

{\footnotesize
Everything reported here is available at \repohref. The entries below are cited
as [R$k$]; they are artifacts rather than literature and are numbered separately
from the bibliography.

\begin{artifactlist}
  \artifact{data}{data/polytopalsimplicialspheres9_4.txt}{Firsching's
    realisations of the $1142$ simplicial types with nine vertices, fetched by
    \artpath{scripts/setup\_data.py}}
  \artifact{firsching}{src/inscribed/enumeration/firsching.py}{reads the data
    file, keeps the inscribed realisations and builds their facets}
  \artifact{certificates}{results/certificates/d4-n09}{the exact certificates
    with $\beta = 468/25$, produced by \artpath{scripts/eliminate.py} and, for
    \texttt{firsching-9-323} and \texttt{firsching-9-324}, by
    \artpath{scripts/eliminate\_non\_inscribable.py}}
  \artifact{verifier}{src/inscribed/certify/verify.py}{the exact verifier of
    \cite[Section~3.3]{medgyes-maxvol-d4}, with
    \artpath{src/inscribed/combinatorics/complex.py} checking coherence}
  \artifact{testnine}{tests/papers/maxvol\_d4\_n9/test\_nine.py}{tests of the
    claims about Firsching's data and the certificates; the checks over all
    $1142$ types run with \texttt{poetry run pytest -m slow}}
\end{artifactlist}
}

\bibliographystyle{amsplain}
\bibliography{references}

\end{document}
